\scanpage{108}

torsor under the module scheme
\[
\left\{\begin{gathered}
\underline{\mathrm{Sc}}_0\simeq\underline K,
\quad\text{where }K\text{ is the kernel of}\\
\text{the exact sequence of }\mathbb Z\text{-modules}\\
\diagram{\begin{tikzcd}[ampersand replacement=\&,column sep=large,row sep=small]
1 \arrow[r,no head] \& K \arrow[r] \&
\deleted{\unsure{\mathbb Z|M}}\mathbb Z^3 \arrow[r] \&
\mathbb Z \arrow[r,no head] \& 0\mathbin{.}\\
\&\& (\alpha_0,\alpha_1,\beta) \arrow[r,mapsto] \&
2\alpha_1-\alpha_0-3\beta \&
\end{tikzcd}}
\end{gathered}\right.
\tag{253}
\]
\margin{$K\simeq\mathbb Z^2$}

On the other hand we have an operation of $\underline M$ on
$\underline{\mathrm{Sc}}$ by \unsure{automorphisms} \unsure{that are inner|that are inner},
which I am going to make explicit; To $\lambda\in\Gamma\underline M$
corresponds the operation
\[
\diagram{\begin{tikzcd}[ampersand replacement=\&,column sep=large]
(\zeta,\eta) \arrow[r,mapsto] \&
\lambda\zeta\rho(\lambda)^{-1},\lambda\eta\sigma(\lambda)^{-1}
\end{tikzcd}}
\]
i.e., additively
\[
\diagram{\begin{tikzcd}[ampersand replacement=\&,column sep=large]
(\zeta,\eta) \arrow[r,mapsto] \&
(\zeta+(1-\rho_M)(\lambda),\eta+(1-\sigma_M)(\lambda))\mathbin{.}
\end{tikzcd}}
\tag{254}
\]

Thus, the operation in question comes from the
$\underline{\mathrm{Sc}}_0$-torsor structure on $\underline{\mathrm{Sc}}$,
and from a homomorphism of \unsure{Groups|groups}
\[
\left\{\diagram{\begin{tikzcd}[ampersand replacement=\&,column sep=large,row sep=large]
\underline M \arrow[r] \& \underline{\mathrm{Sc}}_0\\
\lambda \arrow[r,mapsto] \&
\unsure{(}\underbrace{(1-\rho_M)(\lambda)}_{(c_0+c_1)\lambda_0+(2c_1-c_0)\lambda_1},
\underset{(c_1-c_0)\delta}{(1-\sigma_M)(\lambda)})
\quad\text{where }\lambda=c_0\lambda_0+c_1\lambda_1
\end{tikzcd}}\right.
\tag{255}
\]
which is in fact a homomorphism of \unsure{$M$}-\uline{dules},
associated with the homomorphism of $\mathbb Z$-modules
\[
\left\{\diagram{\begin{tikzcd}[ampersand replacement=\&,column sep=large,row sep=small]
M \arrow[r] \& K \arrow[r,hook] \& \mathbb Z^3\\
c_0\lambda_0+c_1\lambda_1 \arrow[r,mapsto] \&
(\alpha_0=c_0+c_1,\alpha_1=2c_1-c_0,\beta=c_1-c_0) \&
\end{tikzcd}}\right.
\tag{256}
\]

\margin{107}

\scanpage{109}
\mspage{702}

We see at once that this homomorphism is universally injective
(since the composite $M\longrightarrow\mathbb Z^3$ is), hence
\unsure{an} isomorphism for reasons of rank. We have thus found:

\uline{Proposition} The extension
\inserted{$\underline{\mathfrak S}^{+}/\unsure{\mu_{\mathfrak S}}$}
(243) of $(\mathfrak S_3)_{S_0}$ by $\underline M$ is split over $S_0$
--- hence over \unsure{every|every} other base --- the splittings over
$S_0=\operatorname{Spec}\mathbb Z$ correspond \inserted{1-1} to the
splittings of $\mathfrak S^{+}/\mu_{\mathfrak S}$ as an extension of
$\mathfrak S_3$ by $M$. Over an arbitrary base $S$, \deleted{the}
two \deleted{extensions}\inserted{splittings} are conjugate by a unique
section of $\underline M$ over $S$, so the scheme of splittings of
(\unsure{243|247}) is \deleted{\unsure{a}} naturally a torsor
under $\underline M$.

Let us note that it is immediate, over an arbitrary base $S$, that the splittings
of (243) correspond \unsure{one-to-one|one-to-one} to the splittings of
the extension of groups
\[
\diagram{\begin{tikzcd}[ampersand replacement=\&,column sep=large]
1 \arrow[r,no head] \& \underline M(S) \arrow[r,no head] \&
\underline{\mathfrak S}^{+}/\mu_{\mathfrak S}
 (S)^{\Delta}
 \arrow[r] \& \mathfrak S_3 \arrow[r,no head] \& 1
\end{tikzcd}},
\tag{244}
\]
which also identifies with the extension of

\margin{108}

\scanpage{110}

$\mathfrak S_3$ by
$\underline M(S)=\Gamma(\unsure{S},\underline{\mathcal O}_S)
 \otimes_{\mathbb Z}M$, deduced from the universal extension
$\mathfrak S^{+}/\mu_{\mathfrak S}$ of $\mathfrak S_3$ by $M$,
via the homomorphism
\[
\diagram{\begin{tikzcd}[ampersand replacement=\&,column sep=large]
M \arrow[r] \&
\underline M(S)=\unsure{M|\widehat M}\otimes_{\mathbb Z}
\Gamma(S,\underline{\mathcal O}_S)
\end{tikzcd}}.
\]

Let us apply this to the case $S=\operatorname{Spec}(\widehat{\mathbb Z})$,
we find the

\uline{Corollary} The extension $\mathfrak S^{+\wedge}/\mu_{\mathfrak S}$
of $\mathfrak S_3$ by $M^\wedge$ is split, and its splittings form
a torsor under $M^\wedge$ (by automorphisms \unsure{that are inner|that are inner})

Let us now study the splittings of
$\underline G_{\rho,\sigma}/\underline\mu_G$ in its entirety
\unsure{i.e.|:} by \unsure{trying} to determine, for a given splitting
\[
(\rho'=\zeta\rho,\unsure{\sigma'}=\eta\sigma)
\qquad
\left\{\begin{aligned}
\zeta&=a_0\lambda_0+a_1\lambda_1\\
\eta&=b\delta=b(\lambda_1-\lambda_0)
\end{aligned}\right.
\tag{257}
\]
of $\underline{\mathfrak S}^{+}/\underline\mu_{\mathfrak S}$
over $(\mathfrak S_3)_{S_0}$, the liftings
\[
\diagram{\begin{tikzcd}[ampersand replacement=\&,column sep=large]
\mathbb O_{12}^{*} \arrow[r,no head] \&
\underline G_{\rho,\sigma}/\underline\mu_G
\end{tikzcd}}
\]
compatible with the preceding splitting, i.e. which \uline{commute} with
$\rho',\sigma'$. For this

\margin{109}

\scanpage{111}
\mspage{703}

\unsure{I|I} am going to determine the \uline{commutant} of $(\rho',\sigma')$
\deleted{\unsure{of the|in}} in $\underline{G}_{\rho,\sigma}/\unsure{\mu_{\rho}|\mu_G}$,
i.e. the commutant \unsure{of the|of the} partial section subgroup
\[
(\unsure{258})\qquad
\underline{H}(\rho',\sigma')\subset
\unsure{\underline{G}}_{\rho,\sigma}/\unsure{\underline{\mu}_G} .
\]

The image of this commutant in
$(\unsure{\mathfrak{S}}_3)_{S_0}\simeq
(\underline{G}_{\rho,\sigma}/\underline{M})/\mathbb O_{12}^{*}$
is contained in the center of $(\mathfrak{S}_3)_{S_0}$, hence is
\unsure{trivial|trivial}, whence we conclude
\[
\begin{gathered}
(\unsure{259})\qquad
\unsure{\underline{\operatorname{Centr}}}_{\underline{G}_{\rho,\sigma}/\underline{\mu}_G}
(\underline{H}(\rho',\sigma'))\subset
\mathbb O_{12}^{*}\mathbin{.}\underline{M}
\simeq\underline{G}_{\rho,\sigma}^{\unsure{\natural|\sharp}}
\deleted{\unsure{/\mu_G}}\\
\hfill\begin{array}{c}
\Vert\\[-4pt]
\left\{\lambda u_{\underline{\mu}}\ \middle|\
\begin{array}{l}
\lambda\in\Gamma\underline{M}\\
\underline{\mu}\in\Gamma\mathbb O_{12}^{*}
\end{array}\right\}
\end{array}
\end{gathered}
\]
--- where $u_{\underline{\mu}}$ is defined in (150,151). We are therefore
going to determine the centralizer of $\underline{H}(\rho',\sigma')$
by expressing the condition for a $\lambda u_{\underline{\mu}}$
to lie in it, i.e. to commute with $\rho',\sigma'$.

We have\unsure{:}
\[
\begin{aligned}
(\lambda u_{\underline{\mu}})
(\underbrace{\rho'}_{\zeta\rho})
&=\operatorname{int}(\lambda)\,u_{\underline{\mu}}(\zeta\rho)
=\operatorname{int}(\lambda)\,
\underbrace{u_{\underline{\mu}}(\zeta)}_{\zeta^{\mu}}
\underbrace{u_{\underline{\mu}}(\rho)}_{\lambda_1^{\nu}\rho}
=\operatorname{int}(\lambda)(\zeta^{\mu}\lambda_1^{\nu}\rho)\\
&=\lambda\zeta^{\mu}\lambda_1^{\nu}\rho(\unsure{\lambda}^{-1})\rho
\end{aligned}
\]
where $\nu=\nu(\underline{\mu})$, so the relation
$\lambda u_{\underline{\mu}}(\rho')=\rho'$ can also be written
\[
\lambda\zeta^{\mu}\lambda_1^{\nu}\rho(\lambda^{-1})\rho
=\zeta\rho\qquad\text{i.e.}\qquad
\lambda\zeta^{\mu}\lambda_1^{\nu}\rho(\lambda^{-1})=\zeta,
\]
or in additive notation
\[
\lambda+\mu\zeta+\nu\lambda_1-\rho(\lambda)=\zeta,
\qquad\text{i.e.}
\]
taking into account that $\mu-1=2\nu$

\margin{110}

\scanpage{112}

\[
(\unsure{260})\qquad(1-\rho)(\lambda)+2\nu\zeta+\nu\lambda_1=0.
\]

The relation $(\lambda u_{\underline{\mu}})(\sigma')=\sigma'$ becomes explicit, via
\[
\begin{aligned}
(\lambda u_{\underline{\mu}})(\sigma')
&=\operatorname{int}(\lambda)u_{\underline{\mu}}(\eta\sigma)
=\operatorname{int}(\lambda)\left(
\underbrace{u_{\underline{\mu}}(\eta)}_{\eta^{\mu}}
\underbrace{u_{\underline{\mu}}(\sigma)}_{
\substack{\omega^{\nu_2}\sigma=\sigma\\
\text{since }\omega=1\\
\text{in }\underline{G}_{\rho,\sigma}/\mu_G}}
\right)=\lambda\eta^{\mu}\sigma\lambda^{-1}\\
&=\lambda\eta^{\mu}\sigma(\lambda^{-1})\sigma
\end{aligned}
\]
so the relation \inserted{$\lambda u_{\underline{\mu}}(\sigma')=\sigma'$} is written
\[
\lambda\eta^{\mu}\sigma\bigl(\deleted{\unsure{\zeta}}\inserted{\lambda}^{-1}\bigr)\sigma
=\eta\sigma\qquad\text{i.e.}\qquad
\lambda\eta^{\mu}\sigma(\lambda^{-1})=\eta
\]
or again, additively
\[
\lambda+\mu\eta-\sigma(\lambda)=\eta
\]
\[
(\unsure{261|264})\qquad(1-\sigma)\lambda+2\nu\eta=0.
\]

In terms of $a_0,a_1,b$ of (257), these equations are made explicit as
conditions on
\[
\lambda=x\lambda_0+y\lambda_1
\]
by the equations
\[
(\unsure{262|282})\qquad
\left\{
\begin{array}{ll}
-x+y=-2\nu b&\text{(expresses (261))}\\[3pt]
\left.\begin{array}{rcl}
x+y&=&-2\nu a_0\\
-x+2y&=&-\nu(2a_1+1)
\end{array}\right\}&\text{(express (260))}
\end{array}\right.,
\]
where the quantities $a_0,a_1,b$ appearing on the \unsure{right-hand sides}
(\unsure{in|at} addition to the multiplicative factor
$\nu=\nu(\underline{\mu})$) satisfy (250)
\[
(\unsure{263})\qquad-\unsure{a_0}+2a_1-3b+1=0.
\]

The first and third equation (\unsure{262|263}) can be solved for $x,y$
\[
(264)\qquad
\left\{\begin{array}{l}
y=\unsure{\nu|-\nu}(-2a_1+\unsure{2}b-1)\\
x=y+2\nu b=\nu(-2a_1+\unsure{4}b-1)
\end{array}\right.
\]
and the second equation (\unsure{262|261}) becomes, on

\margin{111}

\scanpage{113}
\mspage{704}

substituting these values of $x,y$ into it, a relation linking $a_0,a_1,b,\nu$:
\[
\nu(-4a_1+6b-2)=-2\nu a_0\qquad\text{i.e.}
\]
\[
2\nu(a_0-2a_1+3b-1)=0
\]
which is a consequence of (263)! Thus we find:
\deleted{that for a splitting}

\uuline{Proposition} For \unsure{every} splitting of the extension
$\underline{\mathfrak{S}}^{+}/\mu_{\underline{\mathfrak{S}}}$ of
$(\underline{\mathfrak{S}}_3)_{\unsure{S_0}}$ by $\underline{M}$
(over an arbitrary base), the commutant
\inserted{in $\underline{G}_{\rho,\sigma}/\underline{\mu}_G$}
of the section subgroup is \deleted{\unsure{again|a}} a section
subgroup of \unsure{the} \inserted{\unsure{projection}} the extension
$\underline{G}_{\rho,\sigma}^{\unsure{\natural|\sharp}}$
of $\mathbb O_{12}^{*}$ by $\underline{M}$.
[The lifting
\[
(264)\qquad u_{\rho',\sigma'}:\quad
\diagram{\begin{tikzcd}[ampersand replacement=\&]
\mathbb O_{12}^{*}
\arrow[r] \&
\underline{G}_{\rho,\sigma}^{\unsure{\natural|\sharp}}
\end{tikzcd}}
\]
which is associated with $\rho',\sigma'$ given by (257)
\[
\rho'=\zeta\rho,\quad\unsure{\sigma'}=\eta\sigma\qquad
\begin{array}{l}
\zeta=a_0\lambda_0+a_1\lambda_1,\quad
\eta=b\delta=b(\lambda_1-\lambda_0)\\
a_0,a_1,b\in\Gamma\mathbb{G}_a
\end{array},
\]
is given by
\[
(265)\qquad u_{\rho',\sigma'}(\underline{\mu})
=\lambda u_{\underline{\mu}}
\]
where $u_{\underline{\mu}}$ is defined in (150,151), and where
$\lambda=x\lambda_0+y\lambda_1$, with $x,y\in\Gamma\mathbb{G}_a$
described by 264, \deleted{where} where $\nu=\nu(\underline{\mu})$).

\uline{Corollary} The canonical morphism
\[
(266)\quad
\diagram{\begin{tikzcd}[ampersand replacement=\&,column sep=large]
\underline{\operatorname{Scindages}}\left(
\underline{G}_{\rho,\sigma}/\underline{\mu}_G\,
\deleted{\unsure{\longrightarrow|\text{over}}}
\inserted{\text{over}}\,
\mathbb O_{12}^{*}\times\underline{\mathfrak{S}}_3
\right)\arrow[r]\&
\underline{\operatorname{Scindages}}\left(
\underline{\mathfrak{S}}^{+}\text{ over }\underline{\mathfrak{S}}_3
\right)
\end{tikzcd}}
\]

\margin{112}

\scanpage{114}

is an isomorphism of sheaves \deleted{\unsure{in}}
$(\mathrm{Sch})_{\unsure{S_0}}$, and hence (proposition page 702)
the first is an $\underline{M}$ torsor via conjugation.

This statement is made explicit in a corresponding statement on the set-theoretic
extensions, as in the Corollary to the preceding proposition.
We find in particular, over $S=\operatorname{Spec}(\widehat{\mathbb{Z}})$

\uline{Corollary} The splittings of the extension
$G_{\rho,\sigma}/\mu_G$ of
$\widehat{\mathbb{Z}}^{*}\times\mathfrak{S}_3$ by $\widehat{M}$
form a torsor under $\widehat{M}$, isomorphic
\inserted{(by restriction)} to the torsor of splittings of the extension
$\mathfrak{S}^{+\wedge}/\mu_{\mathfrak{S}}$ \deleted{of} of
$\mathfrak{S}_3$ by $\widehat{M}$.

\figure{\begin{tikzpicture}
\draw (0,0) -- (1.4,0.06);
\end{tikzpicture}}

Variant of the preceding study for the extension $G_{\rho,\sigma}$ of
$\mathbb O_{12}^{*}\times H_{S_0}^{+}$ by
$\underline{M}$ (\uline{without division by}
$\underline{\mu}_G=\{1,\omega\}_{S_0}$). Let us first consider the question of the
splittings of the subextension $\underline{\mathfrak{S}}^{+}$ of
$H_{S_0}^{+}$ by $\underline{M}$ --- \unsure{they|they}
correspond to the pairs

\margin{113}